\section{Background and Related Work}
\label{sec:related}

We organize the related work around the capability ladder introduced above,
tracing the progression from binding through presentation to binary
immunogenicity, and then situating the orthogonal magnitude dimension that our
benchmark targets. We also explicitly delimit a set of tools that produce
continuous outputs but along axes distinct from response magnitude, since these
are the most likely to be mistaken for existing magnitude predictors.

\subsection{The neoantigen prediction ladder}
\label{sec:related:ladder}

Computational prediction of tumour neoantigens has progressed along a capability
ladder: from peptide--MHC binding affinity, to antigen presentation, to binary
T-cell immunogenicity, and---only nominally---to the magnitude of the elicited
T-cell response. Recent surveys frame the discovery pipeline as a sequence of
stages running from HLA typing and variant calling through pMHC presentation to
pMHC recognition \cite{Zhao2025,BiomarkerRes2025}. A recurring lesson across this
literature is that MHC binding is necessary but not sufficient for
immunogenicity: in the TESLA consortium benchmark, only on the order of six
percent of strongly presented candidate peptides were experimentally
immunogenic \cite{TESLA2020}. We emphasize that this ladder describes a
stratification of prediction \emph{targets} rather than a partition of tools into
mutually exclusive bins; modern tools routinely span levels, with binding
predictors emitting both affinity and eluted-ligand scores and immunogenicity
models consuming presentation percentiles as input features.

\subsection{From binding to presentation (L1--L2)}
\label{sec:related:binding}

Binding predictors such as NetMHCpan-4.1 \cite{NetMHCpan2020} and MHCflurry~2.0
\cite{MHCflurry2020} estimate peptide--MHC affinity or eluted-ligand presentation
across thousands of alleles. Presentation-aware models, including PRIME
\cite{PRIME2023} and the eluted-ligand variant of BigMHC \cite{BigMHC2023},
integrate mass-spectrometry eluted-ligand data to estimate the probability that a
peptide is genuinely presented. It is important to recognize that these
continuous outputs quantify binding strength or presentation likelihood, not
response magnitude: the affinity scale is a different physical dimension from the
intensity of a polyclonal T-cell response, and benchmark studies report only weak
correlation between binding scores and ELISpot magnitude
\cite{explorationpub2023,Buckley2022}.

This distinction matters because several established tools produce continuous
scores and could be misread as pre-existing quantitative magnitude predictors.
Three classes warrant explicit delimitation. First, binding-affinity predictors
report a continuous affinity in nanomolar units or a percentile rank
\cite{NetMHCpan2020,MHCflurry2020}. Second, stability predictors such as
NetMHCstabpan estimate a continuous peptide--MHC complex half-life
\cite{NetMHCstabpan2016}. Third, TCR--pMHC models such as pMTnet output a
continuous recognition ranking percentile \cite{pMTnet2021}. None of these
targets the quantity of interest here. Each is continuous along a dimension---
binding affinity, complex stability, or TCR--pMHC recognition propensity---whose
units are not those of a population-level T-cell response measured in
spot-forming cells, and each correlates only weakly with ELISpot magnitude. We
therefore treat these continuous neighbours not as magnitude predictors but as
natural control baselines against which any genuine magnitude regression should
be compared.

\subsection{Immunogenicity classification (L3)}
\label{sec:related:immuno}

The bulk of neoantigen immunogenicity tools target a binary outcome---whether a
given peptide is T-cell immunogenic. Methodologically they cluster into several
families. Classical machine-learning models build on hand-crafted features:
IMPROVE uses a random forest over physicochemical, cleavage, and binding-derived
features \cite{IMPROVE2024}; PredIG uses gradient-boosted trees \cite{PredIG2025};
and ICERFIRE uses a random-forest ensemble \cite{ICERFIRE2024}. Deep
sequence models learn end-to-end representations, as in the convolutional
architecture of DeepImmuno \cite{DeepImmuno2021}, the recurrent model deepHLApan
\cite{deepHLApan2019}, and the eluted-ligand-to-immunogenicity transfer of BigMHC
\cite{BigMHC2023}. Structure-aware models incorporate pMHC surface or
AlphaFold-derived geometry to recover spatial, TCR-facing information lost by
sequence alone, as in NeoaPred \cite{NeoaPred2024} and ImmunoStruct
\cite{ImmunoStruct2025}. A further family pursues multi-modal or cross-domain
integration of presentation, binding, and TCR--pMHC signals \cite{TSCAPE2025},
and earlier tools combined binding with T-cell propensity terms, as in NetTepi
\cite{NetTepi2014} and the integrated pipeline of Seq2Neo \cite{Seq2Neo2022}.
Despite this architectural diversity, every one of these methods ultimately
produces a classification probability or a ranking score rather than a calibrated
quantitative response magnitude.

\subsection{The magnitude gap (L4)}
\label{sec:related:gap}

Within the twelve methods and six reviews we surveyed, no published tool is
explicitly designed to regress the continuous magnitude of the T-cell response and
to report correlation or error against a continuous experimental readout. A recent
review of immunogenicity predictors concludes that the tools it surveys perform
binary classification and that prediction of T-cell response magnitude remains an
unaddressed gap \cite{explorationpub2023}, and a 2025 survey of the field
likewise does not list regression-based magnitude prediction among solved tasks
\cite{Zhao2025}. Existing comparative benchmarks reinforce this boundary by their
scope: prior head-to-head evaluations target binary neoantigen immunogenicity
\cite{Nibeyro2023} or TCR--pMHC binding \cite{Gielis2023}, and recent commentary
attributes the performance ceiling of immunogenicity models to data quality and
dataset bias \cite{OBrien2023}---but none of these benchmarks evaluates the
recovery of continuous response magnitude, which is the gap our benchmark
addresses. The gap persists despite the presence of usable signal. PredIG's
training data carry graded positive-low, intermediate, and high labels that are
binarized before use \cite{PredIG2025}; ICERFIRE states that the labels for a
given peptide--HLA pair are collapsed into a single class
\cite{ICERFIRE2024}; and CNNeoPP's validation ELISpot data distinguish weak from
strong responders, yet the model still emits only a binary call
\cite{CNNeoPP2026}. Across these cases, a magnitude signal that is present in the
underlying data is selectively discarded at the labelling stage.

We are deliberately cautious about how much this gap can be filled. The magnitude
of a polyclonal T-cell response is shaped by donor-specific precursor frequencies
and other host factors that are not observable from peptide and allele sequence
alone, which places a biological ceiling on the variance that any sequence-level
predictor can explain. Accordingly, related work that frames magnitude prediction
as an open frontier must also acknowledge that the achievable correlation is
bounded, and that the practical contribution of continuous prediction lies in
recovering ranking-relevant signal for clinical prioritization rather than in
high-precision regression. Our benchmark is positioned to make exactly this
question empirical: it measures how well existing tools, repurposed as continuous
predictors, track ELISpot magnitude under a harmonized protocol, and it provides
the controls---including the continuous binding, stability, and TCR--pMHC
neighbours delimited above---against which future magnitude predictors should be
judged.
